\originalpreface
基础力学研究物体的位置怎样变化, 以及相互作用怎样改变运动. 本册从物理现象建立向量、导数、力与守恒量, 面向掌握高中代数、三角函数且正在学习单变量微积分的读者. 方程使用国际单位制, 除非题目另有约定. 牛顿理论的基本适用范围是宏观、低速、弱引力情形; 最后一章讨论非惯性系与相对论接口.

主课程为 \href{https://ocw.mit.edu/courses/8-01sc-classical-mechanics-fall-2016/}{MIT 8.01SC Classical Mechanics, 2016 秋}. 教师为 Deepto Chakrabarty、Peter Dourmashkin、Michelle Tomasik、Anna Frebel、Vladan Vuletic. 书面主线为 Peter Dourmashkin 的 \emph{Classical Mechanics: MIT 8.01 Course Notes}, OCW 页面注明 2022 年修订. 没有先修课, 18.01SC 单变量微积分为共修课. 每周两次 2 小时讲课, 一次 1 小时解题课. 第一学年采用 pass/no record. 原大纲列三次期中合计 45\%, 期末 25\%, 作业 10\%, 课堂参与 20\%. 这些比例描述原教学, 本次公开材料没有相应完整考试卷和标准答案, 不补造考试.

\begin{longtable}{@{}p{.13\textwidth}p{.79\textwidth}@{}}
\toprule 周次 & 官方网页课历与本册学习线 \\\midrule
1--3 & 运动学、牛顿定律、圆周运动. 原大纲周次标题还提到动量与冲量, 本册随后独立建立守恒. \\
4--6 & 阻力、约束、连续体系; 动量、冲量、质心; 连续质量流及火箭. \\
7--9 & 动能与功、势能与守恒、碰撞理论及质心参考系. \\
10--12 & 刚体、惯量、力矩、角动量、角冲量、平动与转动及滚动. \\\bottomrule
\end{longtable}
每周作业按原顺序完整纳入, 共 12 套、62 道主问题. 不用网页视频总数冒充已逐字转录视频. 本册冻结书面教材第 1--25、31 章, 共 26 章; 对主要节的实质概念与推导重新组织成下表的连续正文, 而不是保留英文章节顺序的文件堆.
\begin{longtable}{@{}p{.23\textwidth}p{.69\textwidth}@{}}
\toprule 本册正文 & 原教材主要对应章与扩充 \\\midrule
1--2 测量、向量、运动学 & 1--6: 单位、向量、位置速度加速度、二维与圆周运动. \\
3--4 牛顿定律、约束 & 7--9: 力、解题方法、摩擦、绳、滑轮、连续质量与阻力. \\
5--6 动量、质量流 & 10--12: 冲量、守恒、质心、变质量体系. \\
7--8 功、能量、碰撞 & 13--15: 功、势能、能量守恒、碰撞. \\
9--11 刚体、滚动、角动量 & 16--22: 几何、惯量、力矩、静力学、角动量、转动动力学. 8.012 补惯量张量、陀螺与欧拉方程. \\
12 引力与轨道 & 7 的引力与球壳、25: 引力、二体、中心力与开普勒规律. \\
13 振动 & 23--24: 简谐、物理摆、阻尼与受迫运动. \\
14 参考系与接口 & 31 及 8.012 后五讲的加速系、旋转系与相对论基础. \\\bottomrule
\end{longtable}
第 26--30 章的流体与热学属于其他学科, 未纳入本次下载及本册. 第 14 章的热学附录属于跨科内容, 范围单独标识; 机械功、摩擦耗散与内能记账的必要接口仍在正文解释. 本册不重写分析力学的变分法和哈密顿形式, 后续可接仓库分析力学册及 8.03 振动与波.

补充课程为 \href{https://ocw.mit.edu/courses/8-012-physics-i-classical-mechanics-fall-2008/}{MIT 8.012 Physics I: Classical Mechanics, 2008 秋}, 教师 Adam Burgasser. 每周两次 90 分钟讲课、两次 60 分钟习题课. OCW 没有该课完整讲义, 所以不假装能从课历恢复教师逐字讲授. 指定教材为 Daniel Kleppner、Robert J. Kolenkow 的 \emph{An Introduction to Mechanics} (McGraw-Hill, 1973). 原作业里 85 个商业教材题只有书名题号和少量说明, 本册保留这些引用, 不猜造题面. 另有 13 个 MIT 完整题, 均保留解答; 作业 10 可选且不收. 四份公开考试 22 道主问题及全部小问、加分问均纳入, 答案对照四份官方答案. 项目要求及实际运行的数值例也收入同册.

8.012 考核为每周作业 25\%, 课内概念测试 5\%, 项目 5\%, 三次期中各 15\%, 期末 20\%. 考试禁止书籍、笔记、计算器等辅助物, 提供公式表. 课内概念测试没有完整公开题库, 本册不虚构. 以下是其真实课历:
\begin{longtable}{@{}p{.13\textwidth}p{.79\textwidth}@{}}
\toprule 原讲次 & 官方课历主题 \\\midrule
1--5 & 向量; 笛卡尔与极坐标运动; 牛顿定律与力; 牛顿定律应用两讲. 作业 1、2. \\
6--8 & 动量; 第一次碰撞讨论; 火箭运动. 第一次期中, 作业 3. \\
9--12 & 能量; 势能; 能量守恒; 第二次碰撞讨论. 作业 4、5, 第二次期中. \\
13--16 & 定轴转动; 转动惯量; 平动与转动; 刚体运动. 作业 6、7. \\
17--19 & 陀螺两讲; 卡文迪许实验. 作业 8, 第三次期中. \\
20--24 & 惯量张量; 欧拉方程及加速参考系; 加速与旋转参考系; 中心力及开普勒定律; 相对论基础. 作业 9, 项目, 期末. \\\bottomrule
\end{longtable}
官方课历没有列作业 10 的交回时间, 但作业页面公开了它, 故按原件纳入. 不以课历表的缺项删掉已经公开的材料.


正文是按教学依赖重新组织的中文改编, 不是 MIT 官方逐字译本. 97 道完整作业/考试主问题按来源保留原编号与全部已给小问, 教学例题另计. 作业为 AI 独立解答; 考试为官方答案核对后的中文推导, 纠错与额外近似分析明确标识. 图形以 TikZ 重绘; 期末跳水照片作为原题视觉数据保留并逐对象署名记录.

MIT OCW 与作者保留署名; 特别保留原教材第 7 章署名 \copyright\ Peter Dourmashkin 2012. 默认未另标的原件及本册中文改编依 \href{https://creativecommons.org/licenses/by-nc-sa/4.0/}{CC BY-NC-SA 4.0}: 署名、非商业、相同方式共享. 修改包括中文翻译、重编、补充推导、独立题解、SI 定义说明及图形重绘. 明确第三方例外逐对象处理; 书名、题号或参考讨论不使整份混合卷失去 MIT 自编题. 依据 \href{https://mitocw.zendesk.com/hc/en-us/articles/4414756181403-How-is-all-rights-reserved-content-different-from-the-rest-of-OCW-content}{OCW 官方 FAQ}, “all rights reserved” 对象不能自动按默认许可重用; “Courtesy ... Used with permission” 对象须延续专门署名. 本次所取原件未发现明确的前者标记, 不把未提供的商业书题全文补入. 本册不表示 MIT 认可. 来源 URL、文件校验和逐对象记录位于基础力学来源清单.
